% Responsable: A, B, C y D. MAXIMO 1 pagina en total (sin contar referencias).
% Incluir las limitaciones de un modelo lineal alrededor de un punto de operacion.
\section{Marco te\'orico}

% --- Integrante A: modelado de sistemas rotacionales, friccion y frenado
% TODO (A)

% --- Integrante B: identificacion experimental y respuesta al escalon;
%     limitaciones del modelo lineal alrededor de un punto de operacion
La identificaci\'on experimental estima la din\'amica de la planta a
partir de entradas conocidas y de la respuesta medida. Un escal\'on
permite distinguir la magnitud del cambio, la rapidez de la respuesta
y el retardo efectivo. Para el dise\~no de control, esta respuesta puede
aproximarse localmente mediante un modelo de primer orden con
retardo~\cite{skogestad}. En el banco, la fricci\'on y el contacto de la
zapata pueden hacer que la respuesta dependa de la velocidad inicial y
de la amplitud de frenado. Por ello, la aproximaci\'on lineal se eval\'ua
en el intervalo ensayado y con repeticiones independientes.

% --- Integrante C: control PID digital, discretizacion y muestreo
% TODO (C)

% --- Integrante D: frenado realimentado y ABS (alcance del banco didactico)
% TODO (D)
